\documentclass[border=10pt,varwidth=24cm]{standalone}
% deckgl-examples 架构信息图
% 主题：12 个纯本地 deck.gl 参考场景 · 零外部依赖 · 自动化视觉验证
% 读法：从上到下，先看场景家族矩阵，再沿左侧主路径读「数据→渲染→验证」闭环。
% Render: python3 ~/.agents/skills/create-infographic/scripts/render_tikz.py docs/architecture-infographic.tex --workdir . --out-dir docs --engine xelatex --dpi 300 --svg --strict-warnings --strict-geometry --strict-composition --strict-svg --min-labelled-connectors 0.6 --expect-text "deckgl-examples"
\usepackage[UTF8,fontset=none]{ctex}
% 全局轻微字距：消除 dvisvgm 字形墨盒在负 kern 字对上的亚像素重叠（非布局问题）
\defaultfontfeatures{LetterSpace=5}
% CJK 字符间加发丝级 glue：消除跨 text run 的 CJK 字形墨盒亚像素重叠
\xeCJKsetup{CJKglue={\hskip 0.05em plus 0.08em}}
\IfFontExistsTF{Source Han Serif SC}{
  \setmainfont[BoldFont={Source Han Serif SC Bold}]{Source Han Serif SC}
  \setsansfont[BoldFont={Source Han Serif SC Bold}]{Source Han Serif SC}
  \setCJKmainfont[BoldFont={Source Han Serif SC Bold}]{Source Han Serif SC}
  \setCJKsansfont[BoldFont={Source Han Serif SC Bold}]{Source Han Serif SC}
}{
  \PackageWarningNoLine{tikz-infographic}{Source Han Serif SC not found; using Fandol}
  \setCJKmainfont[BoldFont={FandolSong-Bold}]{FandolSong-Regular}
  \setCJKsansfont[BoldFont={FandolHei-Bold}]{FandolHei-Regular}
}
\IfFontExistsTF{Sarasa Mono SC}{
  \setmonofont{Sarasa Mono SC}
  \setCJKmonofont{Sarasa Mono SC}
}{
  \PackageWarningNoLine{tikz-infographic}{Sarasa Mono SC not found; using TeX defaults}
  \setmonofont{Latin Modern Mono}
  \setCJKmonofont{FandolFang-Regular}
}
\usepackage{xcolor}
\usepackage{tikz}
\usepackage{listings}
\usetikzlibrary{arrows.meta,backgrounds,calc,fit,positioning,shapes.geometric,shapes.misc}
\pgfdeclarelayer{edge layer}
\pgfdeclarelayer{background}
\pgfsetlayers{background,edge layer,main}

% ── Semantic palette ──────────────────────────────────────────────────────
\definecolor{Paper}{HTML}{F7F4EE}
\definecolor{Ink}{HTML}{17212B}
\definecolor{Slate}{HTML}{5D6873}
\definecolor{Fog}{HTML}{D9E1E3}
\definecolor{DeepOcean}{HTML}{1B3A5B}
\definecolor{Ocean}{HTML}{2A6F83}
\definecolor{Teal}{HTML}{2A9D8F}
\definecolor{Mint}{HTML}{DDF2EC}
\definecolor{Coral}{HTML}{E76F51}
\definecolor{Gold}{HTML}{E9C46A}
\definecolor{Lavender}{HTML}{A78BFA}
\definecolor{Sky}{HTML}{7BC4D9}
\definecolor{Rose}{HTML}{F4A6A0}
\definecolor{Sand}{HTML}{F0DFB4}

% Scene family colors
\definecolor{FamAgg}{HTML}{2A9D8F}     % 聚合族
\definecolor{FamTemp}{HTML}{E76F51}    % 时序运动族
\definecolor{FamNet}{HTML}{E9C46A}     % 网络流族
\definecolor{FamField}{HTML}{7BC4D9}   % 连续场族
\definecolor{FamPoly}{HTML}{A78BFA}    % 多边形族
\definecolor{FamPoint}{HTML}{F4A6A0}   % 点云族
\definecolor{FamPath}{HTML}{2A6F83}    % 路径拓扑族
\definecolor{FamGlyph}{HTML}{F0DFB4}   % 多变量图元族
\definecolor{FamText}{HTML}{5D6873}    % 标注族

\lstdefinestyle{paper-code}{
  basicstyle=\ttfamily\scriptsize\color{Ink},
  keywordstyle=\bfseries\color{Ocean},
  commentstyle=\itshape\color{Teal!75!black},
  stringstyle=\color{Coral},
  numberstyle=\scriptsize\color{Slate},
  numbers=left,
  numbersep=7pt,
  frame=single,
  rulecolor=\color{Fog},
  backgroundcolor=\color{white},
  breaklines=true,
  columns=fullflexible,
  keepspaces=true,
  showstringspaces=false,
  tabsize=2,
  xleftmargin=8pt,
  framexleftmargin=6pt,
  aboveskip=5pt,
  belowskip=0pt
}

\tikzset{
  every node/.style={font=\rmfamily, inner sep=0pt},
  % Main vertical spine — the rendering data flow
  spine/.style={draw=Fog, line width=5mm, line cap=butt, line join=round},
  stage/.style={
    circle, draw=DeepOcean, fill=DeepOcean, text=white, inner sep=0pt,
    minimum size=7.5mm, font=\rmfamily\scriptsize\bfseries
  },
  stage label/.style={font=\rmfamily\footnotesize\bfseries, text=Ink},
  stage sub/.style={font=\rmfamily\scriptsize, text=Slate},
  % Family token — small color swatch + label
  fam token/.style={
    rounded corners=1.2mm, inner xsep=2mm, inner ysep=1mm,
    font=\rmfamily\scriptsize\bfseries, text=white, align=center
  },
  % Scene card — represents one of the 12 scenes, grouped by family
  scene card/.style={
    draw=Slate!30, fill=white, rounded corners=1.4mm, line width=.5pt,
    text width=25mm, minimum height=16mm,
    font=\rmfamily\scriptsize, text=Ink, align=center,
    inner xsep=2mm, inner ysep=4.5mm
  },
  % Section headings
  section kicker/.style={font=\rmfamily\scriptsize\bfseries, text=Teal},
  section title/.style={font=\rmfamily\large\bfseries, text=Ink},
  section note/.style={font=\rmfamily\scriptsize, text=Slate},
  % Stat tile
  stat/.style={
    rounded corners=2mm, fill=white, draw=Fog, line width=.6pt,
    minimum width=22mm, minimum height=16mm,
    align=center, inner sep=1mm
  },
  stat num/.style={font=\rmfamily\Large\bfseries, text=DeepOcean},
  stat label/.style={font=\rmfamily\scriptsize, text=Slate},
  % Pill/badge
  pill/.style={
    rounded corners=3mm, fill=Mint, text=Ocean,
    inner xsep=2.5mm, inner ysep=1mm,
    font=\rmfamily\scriptsize\bfseries
  },
  pill-warn/.style={
    rounded corners=3mm, fill=Sand, text=Ink,
    inner xsep=2.5mm, inner ysep=1mm,
    font=\rmfamily\scriptsize\bfseries
  },
  pill-coral/.style={
    rounded corners=3mm, fill=Rose!30, text=Coral,
    inner xsep=2.5mm, inner ysep=1mm,
    font=\rmfamily\scriptsize\bfseries
  },
  % Edge styles
  call/.style={
    -{Latex[length=2mm,width=1.4mm]}, draw=DeepOcean,
    line width=.9pt, rounded corners=2mm
  },
  reply/.style={
    -{Latex[length=2mm,width=1.4mm]}, draw=Teal,
    densely dashed, line width=.8pt, rounded corners=2mm
  },
  dataflow/.style={
    -{Stealth[length=1.8mm,width=1.4mm]}, draw=Ocean,
    line width=1pt, rounded corners=1.5mm
  },
  boundary/.style={draw=Slate!40, dash pattern=on 2pt off 2.2pt, line width=.6pt},
  note/.style={font=\rmfamily\scriptsize, text=Slate},
  edge label/.style={font=\rmfamily\scriptsize, text=Slate},
  code chip/.style={
    font=\ttfamily\scriptsize, text=DeepOcean,
    rounded corners=1mm, fill=Fog!40, inner xsep=2mm, inner ysep=1.8mm
  },
  % Package box for dependencies
  pkg box/.style={
    draw=Ocean!40, fill=white, rounded corners=1.8mm, line width=.6pt,
    minimum width=36mm, minimum height=9mm,
    font=\rmfamily\scriptsize\bfseries, text=DeepOcean, align=center
  },
  pkg sub/.style={font=\rmfamily\scriptsize, text=Slate}
}

\begin{document}
\setlength{\parindent}{0pt}
\begin{tikzpicture}[x=1cm,y=1cm]

% ── Paper background ──────────────────────────────────────────────────────
\begin{pgfonlayer}{background}
  \fill[Paper,rounded corners=2.5mm] (-.5,-6.0) rectangle (22.8,29.5);
\end{pgfonlayer}

% ══════════════════════════════════════════════════════════════════════════
% SECTION 1 — 项目总览 (Hero)
% ══════════════════════════════════════════════════════════════════════════
\node[section kicker,anchor=west] at (0,28.9) {DECK.GL EXAMPLES · 架构信息图};
\node[section title,anchor=west,font=\rmfamily\Large\bfseries] at (0,28.15)
  {十二个纯本地 GPU 可视化参考场景};
\node[section note,anchor=west,text width=12cm] at (0,27.5)
  {无需底图、无运行时网络请求；确定性合成数据驱动；从类型检查到画布像素的自动化验证闭环。};

% Stat tiles row
\node[stat] (s12) at (14.2,27.9) {};
\node[stat num] at (s12.center) {12};
\node[stat label,anchor=north,yshift=-1.5mm] at (s12.south) {参考场景};

\node[stat] (s4) at (17.0,27.9) {};
\node[stat num] at (s4.center) {9};
\node[stat label,anchor=north,yshift=-1.5mm] at (s4.south) {视觉家族};

\node[stat] (s0) at (19.6,27.9) {};
\node[stat num,text=Coral] at (s0.center) {0};
\node[stat label,anchor=north,yshift=-1.5mm] at (s0.south) {外部请求};

% Pill row
\node[pill,anchor=west] at (0,26.8) {Vite + TypeScript};
\node[pill,anchor=west] at (2.9,26.8) {deck.gl 9.3};
\node[pill-warn,anchor=west] at (5.4,26.8) {纯合成数据};
\node[pill-coral,anchor=west] at (8.3,26.8) {零底图依赖};
\node[pill,anchor=west] at (10.9,26.8) {透明 PNG 导出};

% Divider
\draw[boundary] (0,26.2) -- (22.3,26.2);

% ══════════════════════════════════════════════════════════════════════════
% SECTION 2 — 技术栈与构建管线
% ══════════════════════════════════════════════════════════════════════════
\node[section kicker,anchor=west] at (0,25.82) {01 · 技术栈};
\node[section title,anchor=west,font=\rmfamily\bfseries] at (0,24.95) {构建与依赖};
\node[section note,anchor=west] at (0,24.4) {单仓库、单入口、四个 deck.gl 子包覆盖全部图层类型。};

% Build pipeline horizontal flow — wider boxes, more vertical room
\node[pkg box, minimum width=30mm] (vite) at (2.0,22.6) {Vite 8.2};
\node[pkg sub,anchor=north,yshift=-1.5mm] at (vite.south) {开发 / 打包};

\node[pkg box, minimum width=30mm] (tsc) at (5.6,22.6) {TypeScript 5.9};
\node[pkg sub,anchor=north,yshift=-1.5mm] at (tsc.south) {类型检查};

\node[pkg box, minimum width=34mm] (deck-core) at (10.0,23.9) {@deck.gl/core};
\node[pkg sub,anchor=north,yshift=-1.5mm] at (deck-core.south) {Deck · 视图系统};

\node[pkg box, minimum width=42mm] (deck-layers) at (14.6,24.8) {@deck.gl/layers};
\node[pkg sub,anchor=north,text width=38mm,align=center,yshift=-1.5mm] at (deck-layers.south) {Arc / Path / Polygon\\Scatterplot / Text};

\node[pkg box, minimum width=42mm] (deck-agg) at (14.6,23.1) {@deck.gl/aggregation-layers};
\node[pkg sub,anchor=north,text width=38mm,align=center,yshift=-1.5mm] at (deck-agg.south) {Contour / Grid\\Heatmap / Hexagon};

\node[pkg box, minimum width=42mm] (deck-geo) at (14.0,21.4) {@deck.gl/geo-layers};
\node[pkg sub,anchor=west,xshift=2mm,yshift=-.7mm] at (deck-geo.east) {TripsLayer};

\node[pkg box, minimum width=28mm] (playwright) at (19.5,23.1) {Playwright};
\node[pkg sub,anchor=north,yshift=-1.5mm] at (playwright.south) {无头验证};

% Build arrows (left side) — route around labels
\begin{pgfonlayer}{edge layer}
  \draw[dataflow] (vite.east) -- (tsc.west);
  \draw[dataflow] (tsc.east) -- (deck-core.west);
  \draw[dataflow] (deck-core.east) -- (12.2,23.9) |- (deck-layers.west);
  \draw[dataflow] (deck-core.east) -- (12.2,23.1) -- (deck-agg.west);
  \draw[dataflow] (deck-core.east) -- (11.8,22.2) -- (11.8,21.4) -- (deck-geo.west);
  \draw[dataflow] (deck-agg.east) -- (playwright.west);
\end{pgfonlayer}
\node[edge label] at (7.55,23.55) {依赖};
\node[edge label] at (17.4,23.4) {渲染输出};

% Command chips sit below the build boxes, inside the paper
\node[code chip] at (2.0,21.3) {npm run build};
\node[code chip] at (5.6,21.3) {npm run validate};

% Divider
\draw[boundary] (0,20.6) -- (22.3,20.6);

% ══════════════════════════════════════════════════════════════════════════
% SECTION 3 — 场景家族矩阵 (core visual)
% ══════════════════════════════════════════════════════════════════════════
\node[section kicker,anchor=west] at (0,20.22) {02 · 场景家族};
\node[section title,anchor=west,font=\rmfamily\bfseries] at (0,19.35) {十二个场景 · 九种视觉家族};
\node[section note,anchor=west,text width=11cm] at (0,18.8)
  {每个场景回答一个典型可视化问题。颜色编码家族身份，同色场景共享底层图层与数据模式。};

% ── Row 1: 聚合族 (2 scenes) ──
\node[fam token, fill=FamAgg, anchor=west] at (0,17.8) {聚合族};

\node[scene card, anchor=west] (hex) at (2.0,17.8) {hexagons\\{\scriptsize\color{Slate}三维密度}};
\node[scene card, anchor=west] (grid) at (5.1,17.8) {grid-cells\\{\scriptsize\color{Slate}规则网格}};

\node[note,anchor=west] at (8.2,17.8) {GPU HexagonLayer / GridLayer};

% ── Row 2: 时序运动族 (1) ──
\node[fam token, fill=FamTemp, anchor=west] at (0,16.6) {时序运动};
\node[scene card, anchor=west] (trips) at (2.0,16.6) {trips\\{\scriptsize\color{Slate}动态轨迹}};
\node[note,anchor=west] at (8.2,16.6) {TripsLayer · 2400 条路径 · 动画帧循环};

% ── Row 3: 网络流 / OD (2) ──
\node[fam token, fill=FamNet, anchor=west] at (0,15.4) {网络流};
\node[scene card, anchor=west] (flows) at (2.0,15.4) {flows\\{\scriptsize\color{Slate}分布式路径}};
\node[scene card, anchor=west] (mig) at (5.1,15.4) {migration-arcs\\{\scriptsize\color{Slate}起止弧}};
\node[note,anchor=west] at (8.2,15.4) {PathLayer + ArcLayer · 节点与流量编码};

% ── Row 4: 连续场 / 标量场 (2) ──
\node[fam token, fill=FamField, anchor=west] at (0,14.2) {连续场};
\node[scene card, anchor=west] (heat) at (2.0,14.2) {heat-islands\\{\scriptsize\color{Slate}热岛叠加}};
\node[scene card, anchor=west] (contour) at (5.1,14.2) {contour-pressure\\{\scriptsize\color{Slate}等值带}};
\node[note,anchor=west] at (8.2,14.2) {HeatmapLayer · ContourLayer · 阈值分带};

% ── Row 5: 多边形 (1) ──
\node[fam token, fill=FamPoly, anchor=west] at (0,13.0) {多边形};
\node[scene card, anchor=west] (parcel) at (2.0,13.0) {parcel-zoning\\{\scriptsize\color{Slate}土地利用}};
\node[note,anchor=west] at (8.2,13.0) {PolygonLayer · 挤压 · 等距视角};

% ── Row 6: 点云 (1) ──
\node[fam token, fill=FamPoint, anchor=west] at (0,11.8) {点云};
\node[scene card, anchor=west] (const) at (2.0,11.8) {constellation\\{\scriptsize\color{Slate}星座场}};
\node[note,anchor=west] at (8.2,11.8) {ScatterplotLayer · 聚类与异常值};

% ── Row 7: 路径拓扑 (1) ──
\node[fam token, fill=FamPath, anchor=west] at (0,10.6) {路径拓扑};
\node[scene card, anchor=west] (river) at (2.0,10.6) {river-network\\{\scriptsize\color{Slate}树状水系}};
\node[note,anchor=west] at (8.2,10.6) {PathLayer · 河流等级驱动宽度};

% ── Row 8: 多变量图元 (1) ──
\node[fam token, fill=FamGlyph, anchor=west] at (0,9.4) {多变量图元};
\node[scene card, anchor=west] (sensor) at (2.0,9.4) {sensor-field\\{\scriptsize\color{Slate}传感器场}};
\node[note,anchor=west] at (8.2,9.4) {大小 + 颜色 + 描边 · 三维编码};

% ── Row 9: 标注 (1) ──
\node[fam token, fill=FamText, anchor=west] at (0,8.2) {标注};
\node[scene card, anchor=west] (label) at (2.0,8.2) {label-atlas\\{\scriptsize\color{Slate}文字图集}};
\node[note,anchor=west] at (8.2,8.2) {TextLayer + ScatterplotLayer};

% Right side: family legend mini
\begin{pgfonlayer}{background}
  \draw[boundary] (12.8,8.0) rectangle (22.3,18.2);
\end{pgfonlayer}
\node[section kicker,anchor=west] at (13.1,17.6) {家族图例};
\node[note,anchor=west,text width=6.5cm] at (13.1,17.0)
  {九个家族对应九种常见的空间可视化问题。同一个家族内的场景共享数据结构、图层组合和交互模式。};

\node[fam token, fill=FamAgg, minimum width=14mm, anchor=west] at (13.1,16.0) {聚合};
\node[note,anchor=west,text width=5cm] at (14.9,16.0) {点数据 → GPU 分箱 → 高度/颜色};

\node[fam token, fill=FamTemp, minimum width=14mm, anchor=west] at (13.1,15.0) {时序};
\node[note,anchor=west,text width=5cm] at (14.9,15.0) {时间戳路径 → 尾迹动画};

\node[fam token, fill=FamNet, minimum width=14mm, anchor=west] at (13.1,14.0) {网络};
\node[note,anchor=west,text width=5cm] at (14.9,14.0) {节点 + 边/弧 → 流量对比};

\node[fam token, fill=FamField, minimum width=14mm, anchor=west] at (13.1,13.0) {场};
\node[note,anchor=west,text width=5cm] at (14.9,13.0) {连续密度 / 等值阈值分带};

\node[fam token, fill=FamPoly, minimum width=14mm, anchor=west] at (13.1,12.0) {多边形};
\node[note,anchor=west,text width=5cm] at (14.9,12.0) {面要素 \kern0.5pt→ 类别 + 挤压};

\node[fam token, fill=FamPoint, minimum width=14mm, anchor=west] at (13.1,11.0) {点云};
\node[note,anchor=west,text width=5cm] at (14.9,11.0) {散点 → 聚类 + 异常};

\node[fam token, fill=FamPath, minimum width=14mm, anchor=west] at (13.1,10.0) {路径};
\node[note,anchor=west,text width=5cm] at (14.9,10.0) {树状拓扑 → 等级宽度};

\node[fam token, fill=FamGlyph, minimum width=14mm, anchor=west] at (13.1,9.0) {图元};
\node[note,anchor=west,text width=5cm] at (14.9,9.0) {多通道视觉编码};

\node[fam token, fill=FamText, minimum width=14mm, anchor=west] at (13.1,8.5) {标注};
\node[note,anchor=west,text width=5cm] at (14.9,8.5) {文本 + 锚点联动};

% Divider
\draw[boundary] (0,7.1) -- (22.3,7.1);

% ══════════════════════════════════════════════════════════════════════════
% SECTION 4 — 数据→渲染 主路径 (left column)  +  验证闭环 (right column)
% ══════════════════════════════════════════════════════════════════════════

% LEFT: data → render flow
\node[section kicker,anchor=west] at (0,6.8) {03 · 渲染管线};
\node[section title,anchor=west,font=\rmfamily\bfseries] at (0,6.15) {数据 → GPU → 画布};
\node[section note,anchor=west,text width=7.5cm] at (0,5.6)
  {确定性伪随机生成器产生可复现数据；deck.gl 在 WebGL 中完成聚合、着色与合成，输出支持透明背景。};

% Vertical spine of 4 stages — centered column
\node[stage] (st1) at (2.0,3.8) {01};
\node[stage label,anchor=west] at (2.9,4.1) {合成数据生成};
\node[stage sub,anchor=west] at (2.9,3.5) {pseudo() · sin 哈希 · 确定性};

\node[stage] (st2) at (2.0,1.8) {02};
\node[stage label,anchor=west] at (2.9,2.1) {图层装配};
\node[stage sub,anchor=west] at (2.9,1.5) {Hexagon / Heatmap / Trips …};

\node[stage] (st3) at (2.0,-0.2) {03};
\node[stage label,anchor=west] at (2.9,0.1) {WebGL 渲染};
\node[stage sub,anchor=west] at (2.9,-0.5) {OrbitView / OrthographicView};

\node[stage] (st4) at (2.0,-2.2) {04};
\node[stage label,anchor=west] at (2.9,-1.9) {透明 PNG 导出};
\node[stage sub,anchor=west] at (2.9,-2.5) {preserveDrawingBuffer · alpha};

\begin{pgfonlayer}{edge layer}
  \draw[spine] (st1.south) -- (st2.north);
  \draw[spine] (st2.south) -- (st3.north);
  \draw[spine] (st3.south) -- (st4.north);
\end{pgfonlayer}
% Spine segment labels — what flows between stages, offset left of the stroke
\node[edge label,anchor=east] at (1.6,2.8) {数据集};
\node[edge label,anchor=east] at (1.6,0.8) {图层栈};
\node[edge label,anchor=east] at (1.6,-1.2) {帧缓冲};

% Data generator code chips — to the RIGHT of stage labels, clear of arrows
\node[code chip,anchor=west] at (7.5,4.3) {makePoints(120k)};
\node[code chip,anchor=west] at (7.5,3.55) {makeTrips(900)};
\node[code chip,anchor=west] at (7.5,2.8) {networkNodes · flows};
\node[code chip,anchor=west] at (7.5,2.05) {branches · parcels · hubs};

% Arrow from chip column into stage 1 — enters from above, label sits above line
\begin{pgfonlayer}{edge layer}
  \draw[call] (7.5,4.5) -- (st1.north|-0,4.5) -- (st1.north);
\end{pgfonlayer}
\node[edge label,anchor=south,yshift=1mm] at (5.5,4.5) {输入数据};

% RIGHT: validation pipeline
\node[section kicker,anchor=west] at (12.8,6.8) {04 · 验证闭环};
\node[section title,anchor=west,font=\rmfamily\bfseries] at (12.8,6.15) {全自动化像素级验证};
\node[section note,anchor=west,text width=7.2cm] at (12.8,5.55)
  {Playwright 启动系统 Chrome，逐场景断言画布内容、控制台零错误、零外网请求。};

% validate.mjs pill — sits below note, above cards
% 局部加字距：消除 'validate' 加粗拉丁字形墨盒的亚像素重叠（非布局问题）
\node[pill-coral,anchor=west] at (12.8,4.9) {\addfontfeature{LetterSpace=12}scripts/validate.mjs};

% Validation stages — cards with generous padding
\node[scene card, fill=white, text width=65mm, minimum height=12mm, align=left, inner xsep=3mm, inner ysep=3.5mm, anchor=west] (v1) at (12.8,4.1)
  {{\bfseries 阻断外网请求}\\{\scriptsize\color{Slate}page.route → 非 127.0.0.1 一律 abort}};

\node[scene card, fill=white, text width=65mm, minimum height=12mm, align=left, inner xsep=3mm, inner ysep=3.5mm, anchor=west] (v2) at (12.8,2.4)
  {{\bfseries 等待场景就绪}\\{\scriptsize\color{Slate}window.\_\_plotDemo.ready === true}};

\node[scene card, fill=white, text width=65mm, minimum height=12mm, align=left, inner xsep=3mm, inner ysep=3.5mm, anchor=west] (v3) at (12.8,0.7)
  {{\bfseries 断言动画帧推进}\\{\scriptsize\color{Slate}trips 场景 frame 必须递增}};

\node[scene card, fill=white, text width=65mm, minimum height=12mm, align=left, inner xsep=3mm, inner ysep=3.5mm, anchor=west] (v4) at (12.8,-1.0)
  {{\bfseries 带样式预览截图}\\{\scriptsize\color{Slate}opaque > 95\% · 彩色像素 > 3500}};

\node[scene card, fill=white, text width=65mm, minimum height=12mm, align=left, inner xsep=3mm, inner ysep=3.5mm, anchor=west] (v5) at (12.8,-2.7)
  {{\bfseries 透明 PNG 导出校验}\\{\scriptsize\color{Slate}alpha 透明 > 10\% · 内容像素 > 2500}};

% Validation flow arrows
\begin{pgfonlayer}{edge layer}
  \draw[dataflow] (v1.south) -- (v2.north);
  \draw[dataflow] (v2.south) -- (v3.north);
  \draw[dataflow] (v3.south) -- (v4.north);
  \draw[dataflow] (v4.south) -- (v5.north);
\end{pgfonlayer}
% Edge labels in the gaps, beside the stroke
\node[edge label,anchor=west] at (16.75,3.25) {请求白名单};
\node[edge label,anchor=west] at (16.75,1.55) {ready 信号};
\node[edge label,anchor=west] at (16.75,-0.15) {递增帧};
\node[edge label,anchor=west] at (16.75,-1.85) {预览 PNG};

% Bottom: outputs
\draw[boundary] (0,-4.0) -- (22.3,-4.0);

\node[section kicker,anchor=west] at (0,-4.4) {产物};
\node[code chip,anchor=west] at (0,-5.25) {out/*-transparent.png};
\node[code chip,anchor=west] at (4.4,-5.25) {out/*-preview.png};
\node[code chip,anchor=west] at (8.2,-5.25) {catalog.json};
\node[code chip,anchor=west] at (10.7,-5.25) {dist/};

\node[note,anchor=east] at (22.3,-5.25) {deckgl-examples · v0.1.0};

\end{tikzpicture}
\end{document}
